% ============ 7 X_et 上的层范畴 (原书 pp.50-57) ============
\chapter{\texorpdfstring{$X_{\mathrm{et}}$ 上的层范畴}{X\_et 上的层范畴}}\label{sec:7}
\index{Abel 范畴}
\index{预层范畴}
\index{层范畴}
\index{伴随函子}
\index{内射对象}
\index{局部满射}
\index{层化}
\index{Kummer 序列}
\index{Artin-Schreier 序列}
\index{代数空间}
\index{Noether 景}

本节研究 $X_{\mathrm{et}}$ 上 Abel 群层的范畴.特别地，我们证明它是 Abel 范畴.

\subsection*{范畴的一般知识}

设 $T$ 是加性范畴---回忆这意味着集合 $\operatorname{Hom}(A, B)$ 带有 Abel 群结构，使得复合映射双加性，且 $T$ 中每个有限对象族都有直和.$T$ 中的序列
\[
0 \to A \to B \xrightarrow{\alpha} C
\]
正合，如果对所有对象 $T \in T$，
\[
0 \to \operatorname{Hom}(T, A) \to \operatorname{Hom}(T, B) \to \operatorname{Hom}(T, C)
\]
正合，这时 $A$ 称为 $\alpha$ 的核.序列
\[
A \xrightarrow{\beta} B \to C \to 0
\]
正合，如果对所有对象 $T \in T$，
\[
0 \to \operatorname{Hom}(C, T) \to \operatorname{Hom}(B, T) \to \operatorname{Hom}(A, T)
\]
正合，这时 $C$ 称为 $\beta$ 的余核.

设 $T$ 是每个态射都有核与余核的加性范畴.设 $\alpha \colon A \to B$ 是 $T$ 中的态射.$\alpha$ 的余核的核称为 $\alpha$ 的像，$\alpha$ 的核的余核称为 $\alpha$ 的余像.有从 $\alpha$ 的余像到 $\alpha$ 的像的典范态射，若它总是同构，则 $T$ 称为 Abel 范畴.

设 $T$ 是 Abel 范畴.把两个定义结合起来，得到短正合序列
\[
0 \to A \to B \to C \to 0
\]
的概念，从而得到长正合序列的概念.

从一个加性范畴到另一个加性范畴的函子，我们总指加性函子，即要求映射 $\operatorname{Hom}(X, Y) \to \operatorname{Hom}(FX, FY)$ 是 Abel 群同态.

现在设 $C$ 是小范畴\footnote{范畴称为小的，如果它的对象构成集合(而不是类).}，$T$ 是从 $C$ 到 Abel 群范畴 $\mathbf{Ab}$ 的反变函子 $F \colon C \to \mathbf{Ab}$ 的范畴.态射 $\alpha \colon F_1 \to F_2$ 是自然变换.于是，$\alpha$ 给 $C$ 的每个对象 $U$ 配上同态 $\alpha(U) \colon F_1(U) \to F_2(U)$，使得对每个态射 $\varphi \colon V \to U$，图
\[
\begin{tikzcd}[row sep=small, column sep=small]
  F_1(U) \arrow[r, "\alpha(U)"] \arrow[d, "F_1(\varphi)"'] & F_2(U) \arrow[d, "F_2(\varphi)"] \\
  F_1(V) \arrow[r, "\alpha(V)"'] & F_2(V)
\end{tikzcd}
\]
交换.

显然，$T$ 在 $\operatorname{Hom}(F_1, F_2)$ 的显然群结构和直和
\[
(\oplus_i F_i)(U) = \oplus_i F_i(U)
\]
下成为加性范畴.

我们把验证如下论断留作(简单的)习题：

\begin{quote}
\refstepcounter{theorem}\label{item:7.1}
\textbf{7.1} (a) $C$ 上反变函子的序列 $\cdots \to F_i \to G_i \to H_i \to \cdots$ 正合，当且仅当对所有 $U \in C$，序列 $\cdots \to F_i(U) \to G_i(U) \to H_i(U) \to \cdots$ 正合.

(b) $C$ 上反变函子的范畴是 Abel 的.
\end{quote}

\begin{remark}\label{rem:7.2}
在范畴的定义中，要求对每对对象 $X, Y$，$\operatorname{Hom}(X, Y)$ 是集合(而不是类).若 $C$ 不小，则 $C$ 上两个函子之间的自然变换 $\alpha \colon F \to G$ 可能不构成集合.
\end{remark}

\textbf{伴随函子} 设 $C_1$ 与 $C_2$ 是 Abel 范畴.函子 $L \colon C_1 \to C_2$ 与 $R \colon C_2 \to C_1$ 称为伴随的，如果
\[
\operatorname{Hom}_{C_2}(LX_1, X_2) \cong \operatorname{Hom}_{C_1}(X_1, RX_2),
\]
关于 $X_1$ 与 $X_2$ 函子性地成立.也说 $L$ 是 $R$ 的左伴随，$R$ 是 $L$ 的右伴随.若 $L$ 有右伴随，则右伴随唯一(到唯一的同构).

回忆加性范畴的对象 $I$ 称为内射的，如果 $\operatorname{Hom}(\cdot, I)$ 是正合函子.

我们把验证如下论断留作(简单的)习题：

\begin{quote}
\refstepcounter{theorem}\label{item:7.3}
\textbf{7.3} (a) 有左伴随的函子 $R$ 左正合(事实上保持乘积与逆向极限).

(b) 有右伴随的函子 $L$ 右正合(事实上保持直和与正向极限).

(c) 有正合左伴随的函子 $R$ 保持内射对象.
\end{quote}

\subsection*{预层范畴}

设 $X$ 是簇(或概形).按定义，$X_{\mathrm{et}}$ 上 Abel 群的预层就是反变函子 $\mathbf{Et}/X \to \mathbf{Ab}$.为了把上一节的讨论应用于它们，严格地说，我们应该把 $\mathbf{Et}/X$ 换成小范畴.例如，可以把 $\mathbf{Et}/X$ 换成如下形式的平展映射 $U \to X$ 的集合：$U$ 由如下形状的环的商所联系的簇(或概形)拼接而成：$A[T_1, T_2, \ldots]$，其中 $A = \Gamma(V, \mathcal{O}_X)$($V$ 是 $X$ 的某个开仿射子集)，$\{T_1, T_2, \ldots\}$ 是固定符号集.我们将忽略这个问题，并把 $X_{\mathrm{et}}$ 上 Abel 群预层的范畴 $\mathrm{PreSh}(X_{\mathrm{et}})$ 当作 $\mathbf{Et}/X$ 上所有反变函子的范畴来谈论.它是 Abel 范畴，其中 $P' \to P \to P''$ 正合，当且仅当对一切平展映射 $U \to X$，$P'(U) \to P(U) \to P''(U)$ 正合.核，余核，乘积，直和，逆向极限，正向极限等在 $\mathrm{PreSh}(X_{\mathrm{et}})$ 中以显然的方式形成：对每个平展映射 $U \to X$ 构造核，余核等，并取诱导的限制映射.

\subsection*{层范畴}

设 $X$ 是簇(或概形).我们定义范畴 $\mathrm{Sh}(X_{\mathrm{et}})$ 为 $\mathrm{PreSh}(X_{\mathrm{et}})$ 的以 $X_{\mathrm{et}}$ 上 Abel 群层为对象的全子范畴.于是，给出 $X_{\mathrm{et}}$ 上层 $F' \to F$ 的态射就是给出函子的自然变换 $F' \to F$.显然，$\mathrm{Sh}(X_{\mathrm{et}})$ 是加性范畴.我们考察 $\mathrm{Sh}(X_{\mathrm{et}})$ 中的正合性意味着什么.

层(或预层)的态射 $\alpha \colon F \to F'$ 称为局部满的，如果对每个 $U$ 与 $s \in F'(U)$，存在覆盖 $(U_i \to U)$，使得对每个 $i$，$s|_{U_i}$ 在 $F(U_i) \to F'(U_i)$ 的像中.

\begin{lemma}\label{lem:7.4}
设 $\alpha \colon F \to F'$ 是 $X_{\mathrm{et}}$ 上层之间的同态.如下论断等价：
\begin{enumerate}
  \item 层序列 $F \xrightarrow{\alpha} F' \to 0$ 正合；
  \item 映射 $\alpha$ 局部满；
  \item 对每个几何点 $\bar{x} \to X$，映射 $\alpha_{\bar{x}} \colon F_{\bar{x}} \to F'_{\bar{x}}$ 满射.
\end{enumerate}
\end{lemma}

\begin{proof}
(b) $\Rightarrow$ (a).设 $\beta \colon F' \to T$ 是层映射，使得 $\beta \circ \alpha = 0$；我们必须证明这蕴含 $\beta = 0$.

设 $s' \in F'(U)$($U \to X$ 平展).按假设，存在覆盖 $(U_i \to U)_{i \in I}$ 与 $s_i \in F(U_i)$，使得 $\alpha(s_i) = s'|_{U_i}$.现在
\[
\beta(s')|_{U_i} = \beta(s'|_{U_i}) = \beta \circ \alpha(s_i) = 0, \qquad \text{对所有 } i.
\]
由于 $T$ 是层，这蕴含 $\beta(s') = 0$.

(a) $\Rightarrow$ (c).设对某个 $\bar{x} \in X$，$\alpha_{\bar{x}}$ 不满，并设 $\Lambda \neq 0$ 是 $F_{\bar{x}} \to F'_{\bar{x}}$ 的余核.设 $\Lambda_x$ 是上一节定义的层；于是对 $X_{\mathrm{et}}$ 上任意层 $G$，$\operatorname{Hom}(G, \Lambda_x) = \operatorname{Hom}(G_{\bar{x}}, \Lambda)$.映射 $F'_{\bar{x}} \to \Lambda$ 定义非零态射 $F' \to \Lambda_x$，它与 $F \to F'$ 的复合为零(因为它对应于复合 $F_{\bar{x}} \to F'_{\bar{x}} \to \Lambda$).因此，$F \xrightarrow{\alpha} F' \to 0$ 不正合.

(c) $\Rightarrow$ (b).设 $U \to X$ 平展，$\bar{u} \to U$ 是 $U$ 的几何点.复合 $\bar{u} \to U \to X$ 是几何点；记它为 $\bar{x}$.$\bar{u}$ 的平展邻域通过与 $U \to X$ 复合给出 $\bar{x}$ 的平展邻域，而以这种方式产生的 $\bar{x}$ 的平展邻域是共尾的；因此对 $X_{\mathrm{et}}$ 上每个层 $F$，$F_{\bar{u}} \cong F_{\bar{x}}$.

于是，假设蕴含对 $U$ 的每个几何点 $\bar{u} \to U$，$F_{\bar{u}} \to F'_{\bar{u}}$ 满射.设 $s \in F'(U)$.设 $u \in U$，$\bar{u} \to U$ 是像为 $u$ 的 $U$ 的几何点.因为 $F_{\bar{u}} \to F'_{\bar{u}}$ 满射，存在其像包含 $u$ 的平展映射 $V \to U$，使得 $s|_{V}$ 在 $F(V) \to F'(V)$ 的像中.对足够多的 $u \in U$ 应用这个论断，得到所求的覆盖.
\end{proof}

\begin{proposition}\label{prop:7.5}
设
\[
0 \to F' \to F \to F''
\]
是 $X_{\mathrm{et}}$ 上的层序列.如下论断等价：
\begin{enumerate}
  \item 该序列在层范畴中正合；
  \item 对所有平展映射 $U \to X$，序列
  \[
  0 \to F'(U) \to F(U) \to F''(U)
  \]
  正合；
  \item 对 $X$ 的每个几何点 $\bar{x} \to X$，序列
  \[
  0 \to F'_{\bar{x}} \to F_{\bar{x}} \to F''_{\bar{x}}
  \]
  正合.
\end{enumerate}
\end{proposition}

\begin{proof}
在下一小节中，我们证明函子 $i \colon \mathrm{Sh}(X_{\mathrm{et}}) \to \mathrm{PreSh}(X_{\mathrm{et}})$ 有左伴随 $a$.因此(见~\ref{item:7.3})，$i$ 左正合，这证明了 (a) 与 (b) 的等价.

Abel 群正合序列的正向极限正合，从而 (b) 蕴含 (c).逆命题不难，利用 $s \in F(U)$ 为零，当且仅当对 $U$ 的所有几何点 $\bar{u}$, $s_{\bar{u}} = 0$(比较前一个命题 (c) 蕴含 (b) 的证明).
\end{proof}

\begin{proposition}\label{prop:7.6}
设
\[
0 \to F' \to F \to F'' \to 0
\]
是 $X_{\mathrm{et}}$ 上 Abel 群层的序列.如下论断等价：
\begin{enumerate}
  \item 该序列在层范畴中正合；
  \item 映射 $F \to F''$ 局部满，且对所有开平展映射 $U \to X$，序列
  \[
  0 \to F'(U) \to F(U) \to F''(U)
  \]
  正合；
  \item 对每个几何点 $\bar{x} \to X$，序列
  \[
  0 \to F'_{\bar{x}} \to F_{\bar{x}} \to F''_{\bar{x}} \to 0
  \]
  正合.
\end{enumerate}
\end{proposition}

\begin{proof}
结合最后两个命题.
\end{proof}

\begin{remark}\label{rem:7.7}
在最后三个命题的 (c) 中，只需对簇 $X$ 的每个闭点 $x$ 的一个几何点 $\bar{x} \to x \in X$，或对概形 $X$ 的每个点 $x$ 的一个几何点 $\bar{x} \to x \in X$ 检验条件.
\end{remark}

\begin{proposition}\label{prop:7.8}
$X_{\mathrm{et}}$ 上 Abel 群层的范畴 $\mathrm{Sh}(X_{\mathrm{et}})$ 是 Abel 范畴.
\end{proposition}

\begin{proof}
从态射的余像到其像的映射是同构，因为它在茎上同构.
\end{proof}

\begin{example}\label{ex:7.9}
(a) (Kummer 序列).设 $n$ 是不被 $X$ 的任何剩余域的特征整除的整数.例如，若 $X$ 是特征 $p \neq 0$ 的域 $k$ 上的簇，则要求 $p$ 不整除 $n$.考虑序列
\[
0 \to \mu_n \to \mathbb{G}_{m} \xrightarrow{t \mapsto t^{n}} \mathbb{G}_{m} \to 0.
\]
由命题~\ref{prop:7.6} 与例~\ref{ex:6.13}(b)，为证明它正合，只需验证对 $X$ 的每个严格局部环 $A = \mathcal{O}_{X, \bar{x}}$，
\[
0 \to \mu_n(A) \to A^{\times} \xrightarrow{t \mapsto t^{n}} A^{\times} \to 0
\]
正合.除第二个 $A^{\times}$ 处外这都是显然的，而在这里我们必须证明 $A^{\times}$ 的每个元素都是 $n$ 次幂.但 $\frac{d(T^{n} - a)}{dT} = nT^{n-1} \neq 0$(在 $A$ 的剩余域中)，因此 $T^{n} - a$ 在 $A[T]$ 中分裂.

(b) (Artin-Schreier 序列).设 $X$ 是特征 $p \neq 0$ 的域 $k$ 上的簇，考虑序列
\[
0 \to \mathbb{Z}/p\mathbb{Z} \to \mathbb{G}_{a} \xrightarrow{t \mapsto t^{p} - t} \mathbb{G}_{a} \to 0.
\]
同样，为证明该序列正合，必须验证对 $X$ 的每个严格局部环 $A = \mathcal{O}_{X, \bar{x}}$，
\[
0 \to \mathbb{Z}/p\mathbb{Z} \to A \xrightarrow{t \mapsto t^{p} - t} A \to 0
\]
正合.除第二个 $A$ 处外这都是显然的，但 $\frac{d(T^{p} - T - a)}{dT} = -1 \neq 0$(在 $A$ 的剩余域中)，因此 $T^{p} - T - a$ 在 $A[T]$ 中分裂.
\end{example}

\begin{remark}\label{rem:7.10}
若 $p$ 整除 $X$ 的某个剩余域的特征，则
\[
0 \to \mu_p \to \mathbb{G}_{m} \xrightarrow{t \mapsto t^{p}} \mathbb{G}_{m} \to 0
\]
对 $X$ 上的平展拓扑不正合.然而，它对平坦拓扑正合，因为对任何 $a \in \Gamma(U, \mathcal{O}_U)$($U$ 仿射)，方程 $T^{p} - a$ 定义 $U$ 的平坦覆盖.因此 $\mathbb{G}_{m}(U)$ 的每个元素对平坦拓扑局部地都是 $p$ 次幂.
\end{remark}

\subsection*{与预层相伴的层}

本小节中，``层''(或``预层'')指集合的层(或预层).

\begin{definition}\label{def:7.11}
设 $P \to aP$ 是从预层 $P$ 到(某个景上的)层的同态；如果从 $P$ 到层的所有其他同态都唯一穿过 $P \to aP$，即对所有层 $F$，
\[
\operatorname{Hom}(P, F) \cong \operatorname{Hom}(aP, F),
\]
则 $aP$ 称为与 $P$ 相伴的层(或 $P$ 的层化).
\end{definition}

显然，带映射 $P \to aP$ 的 $aP$ 若存在，则唯一到唯一的同构.本小节余下部分中，我们在各种情形中构造 $aP$.首先给出 $F$ 是与 $P$ 相伴的层的判别准则.

设 $P$ 是预层.截面 $s_1, s_2 \in P(U)$ 称为局部相等的，如果对 $U$ 的某个覆盖 $(U_i \to U)_{i \in I}$ 中的所有 $U_i$，$s_1|_{U_i} = s_2|_{U_i}$.层判别准则蕴含层的局部相等截面相等.

\begin{lemma}\label{lem:7.12}
设 $i \colon P \to F$ 是从预层 $P$ 到层 $F$ 的同态.假定：
\begin{enumerate}
  \item $P$ 中在 $F(U)$ 中有相同像的截面只有那些局部相等的截面，且
  \item $i$ 局部满.
\end{enumerate}
则 $(F, i)$ 是与 $P$ 相伴的层.
\end{lemma}

\begin{proof}
设 $i' \colon P \to F'$ 是从 $P$ 到层的第二个映射，对某个 $U$，设 $s \in F(U)$.由 (b)，知道对 $U$ 的某个覆盖 $(U_i \to U)$，存在 $s_i \in P(U_i)$，使得 $i(s_i) = s|_{U_i}$.由 (a)，$i'(s_i) \in F'(U_i)$ 不依赖于 $s_i$ 的选取，而且 $i'(s_i)$ 与 $i'(s_j)$ 在 $F'(U_i \times_{U} U_j)$ 上的限制一致.我们定义 $\alpha(s)$ 为 $F'(U)$ 中限制为所有 $i$ 的 $i'(s_i)$ 的唯一元素.则 $s \mapsto \alpha(s) \colon F \to F'$ 是满足 $i \circ \alpha = i'$ 的唯一同态.
\end{proof}

\begin{lemma}\label{lem:7.13}
设 $P$ 是层 $F$ 的子预层.对每个 $U$，设 $P'(U)$ 是 $F(U)$ 中如下元素的集合：存在覆盖 $(U_i \to U)_{i \in I}$，使得对每个 $i$，$s|_{U_i} \in P(U_i)$(即 $s$ 局部在 $P$ 中).则 $P'$ 是 $F$ 的子层，且 $P \to P'$ 局部满.
\end{lemma}

\begin{proof}
证明是平凡的.
\end{proof}

我们把 $P'$ 称为 $F$ 的由 $P$ 生成的子层.

要构造 $aP$，只需构造满足 (7.12a) 的层 $F$ 与同态 $P \to F$，因为这时 $F$ 的由 $P$ 的像生成的子层就是与 $P$ 相伴的层.

现在我们处理 $X_{\mathrm{et}}$ 上的层.

\begin{lemma}\label{lem:7.14}
若 $i \colon P \to F$ 满足引理~\ref{lem:7.12} 的条件 (a) 和 (b)，则对所有几何点 $\bar{x}$，
\[
i_{\bar{x}} \colon P_{\bar{x}} \to F_{\bar{x}}
\]
是同构.
\end{lemma}

\begin{proof}
这容易由各个定义推出.
\end{proof}

对每个 $x \in X$，选取像为 $x$ 的几何点 $\bar{x} \to X$.对 $X$ 上的预层 $P$，定义 $P^{*} = \prod (P_{\bar{x}})_{\bar{x}}$，其中 $(P_{\bar{x}})_{\bar{x}}$ 是按 \S 6 的方式与 Abel 群 $P_{\bar{x}}$ 相伴的层.则 $P^{*}$ 是层，且自然映射 $P \to P^{*}$ 满足 (\ref{lem:7.12}) 的条件 (a).

\begin{theorem}\label{thm:7.15}
对 $X_{\mathrm{et}}$ 上的每个预层 $P$，存在相伴层 $i \colon P \to aP$.映射 $i$ 在茎上诱导同构 $P_{\bar{x}} \to (aP)_{\bar{x}}$.函子 $a \colon \mathrm{PreSh}(X_{\mathrm{et}}) \to \mathrm{Sh}(X_{\mathrm{et}})$ 正合.
\end{theorem}

\begin{proof}
取 $aP$ 为 $P^{*}$ 的由 $i(P)$ 生成的子层.则 $i \colon P \to aP$ 满足引理~\ref{lem:7.12} 的条件 (a) 和 (b)，前两个论断由它推出.若 $P' \to P \to P''$ 是 Abel 群的正合序列，则对所有 $x \in X$，$P'_{\bar{x}} \to P_{\bar{x}} \to P''_{\bar{x}}$ 正合(因为 Abel 群正合序列的正向极限正合).但最后的序列可以等同于 $(aP')_{\bar{x}} \to (aP)_{\bar{x}} \to (aP'')_{\bar{x}}$，这表明 $aP' \to aP \to aP''$ 正合.
\end{proof}

\begin{example}\label{ex:7.16}
群 $\Lambda$ 定义常值预层 $P_{\Lambda}$，使得对一切 $U \neq \varnothing$，$P_{\Lambda}(U) = \Lambda$.与 $P_{\Lambda}$ 相伴的层是 $F_{\Lambda}$(见 \S 6).

若 $P$ 是 Abel 群层，则 $aP$ 也是.
\end{example}

\begin{remark}\label{rem:7.17}
设 $(F_i)_{i \in I}$ 是 $X_{\mathrm{et}}$ 上 Abel 群层的族.

设 $P$ 是如下预层：对所有平展映射 $U \to X$，$P(U) = \prod_{i \in I} F_i(U)$，限制映射显然.则 $P$ 是层，且它是诸 $F_i$ 的乘积.类似的评注适用于逆向极限与核.

设 $P$ 是如下预层：对所有平展映射 $U \to X$，$P(U) = \bigoplus_{i \in I} F_i(U)$，限制映射显然.一般地，$P$ 不是层，但 $aP$ 是 $\mathrm{Sh}(X_{\mathrm{et}})$ 中诸 $F_i$ 的直和.类似的评注适用于正向极限与余核.

为免读者以为层范畴``就像 Abel 群范畴一样''，让我指出：一般地它没有足够的投射对象；层的正合序列的乘积也不一定正合.
\end{remark}

\begin{aside}\label{asd:7.18}
设 $P$ 是拓扑空间 $X$ 上的集合预层.定义 $E(P)$ 为 $\coprod_{x \in X} P_x$，设 $\pi \colon E(P) \to X$ 是把 $P_x$ 的每个元素送到 $x$ 的映射.$P(U)$ 的 $s$ 定义 $\pi$ 在 $U$ 上的截面 $s \colon U \to E(P)$，即 $x \mapsto s_x$.赋予 $E(P)$ 使所有映射 $s \colon U \to E(P)$ 连续的最细拓扑：于是 $V \subset E(P)$ 开，当且仅当对所有开集 $U \subset X$ 与 $s \in P(U)$，$s^{-1}(V)$ 在 $U$ 中开.

[邀请读者画一幅图：$X$ 表示成水平线，每个茎 $P_x$ 表示成 $x$ 上方的一条短竖线；``层'',``茎'',``截面''这些术语的由来现在应当清楚了.]

对 $X$ 的每个开子集 $U$，设 $F(U)$ 是 $\pi$ 在 $U$ 上的一切连续截面 $s \colon U \to E(P)$ 的集合.显然的映射 $P \to F$ 把 $F$ 实现为与 $P$ 相伴的层.[空间 $E(P)$ 是与 $P$ 相伴的``\'etal\'e 空间''---见 Godement, R., \emph{Th\'eorie des Faisceaux}, Hermann, 1964, II 1.2.可以避免使用这些空间---事实上，Grothendieck 已把它们从数学中驱逐出去了---但它们相当有用，例如用于定义层的逆像.有趣的是，对平展拓扑有一个(需要代数空间的)类比，它例如对定义 Frobenius 映射在上同调群上的作用是有用的.这一方面在讲义的这个版本中强调得不够.
\end{aside}

\begin{aside}\label{asd:7.19}
景 $T$ 上的集合预层 $P$ 称为分离的，如果两个截面只要局部相等就相等.对预层 $P$，定义 $P^{+}$ 为如下预层：
\[
P^{+}(U) = \varinjlim \operatorname{Ker}\left( \prod P(U_i) \rightrightarrows \prod P(U_i \times_{X} U_j) \right),
\]
---极限遍历所有覆盖(假定极限存在).不难证明：
\begin{enumerate}
  \item 若 $P$ 分离，则 $P^{+}$ 是层；
  \item 对任意 $P$，$P^{+}$ 分离.
\end{enumerate}
于是，对每个预层 $P$，$P^{++}$ 是层.事实上，它就是与 $P$ 相伴的层.

正如 Waterhouse 指出的(Pac. J. Math 57 (1975), 597--610)，存在极限不存在的景：问题是极限遍历的是一个类而非集合，不能换成集合上的极限.除了可能 fpqc 景外，这似乎对实际使用的任何景都不是问题.
\end{aside}

\subsection*{Artin 代数空间}

Zariski 拓扑允许我们把仿射簇拼接起来得到一般簇.可以问平展拓扑是否允许我们把仿射簇拼接起来得到更一般的对象.答案是肯定的！

为简单起见，我在固定的代数闭域 $k$ 上工作，只考虑 $k$-簇.

赋予 $\mathbf{Aff}/k$(所有仿射 $k$-簇的范畴)平展拓扑.该景上的层是满足如下条件的反变函子 $A \colon \mathbf{Aff}/k \to \mathbf{Sets}$：对 $\mathbf{Aff}/k$ 中平展映射的每个满族 $(U_i \to U)_{i \in I}$ 满足层条件.注意 $A$ 也可以看作从仿射 $k$-代数到 $\mathbf{Sets}$ 的共变函子.对每个 $k$-簇 $V$，把仿射 $k$-代数 $A$ 送到 $V$ 中坐标在 $A$ 里的点的集合的函子 $A \mapsto V(A)$ 是 $(\mathbf{Aff}/k)_{\mathrm{et}}$ 上的层；此外，由 Yoneda 引理，$k$-簇的范畴可以看作 $(\mathbf{Aff}/k)_{\mathrm{et}}$ 上层范畴的全子范畴.

设 $U$ 是仿射 $k$-簇.$U$ 上的平展等价关系是子簇 $R \subset U \times U$，使得：
\begin{enumerate}
  \item 对每个仿射 $k$-代数 $A$，$R(A)$ 是 $U(A)$ 上的等价关系(按通常意义---这些是集合)；
  \item 包含 $R \hookrightarrow U \times U$ 与投影映射 $U \times U \rightrightarrows U$ 的复合满且平展.
\end{enumerate}
$(\mathbf{Aff}/k)_{\mathrm{et}}$ 上的层 $A$ 是代数空间，如果存在仿射 $k$-簇 $U$，$U$ 上的平展等价关系 $R$ 和把 $A$ 实现为 $U$ 除以 $R$ 的层论商层的层映射 $U \to A$.最后这个条件意味着，
\begin{enumerate}
  \item 对每个仿射 $k$-代数 $A$ 与 $s_1, s_2 \in U(A)$，$s_1$ 与 $s_2$ 在 $A(A)$ 中有相同的像，当且仅当 $(s_1, s_2) \in R(A)$；
  \item 层映射 $U \to A$ 局部满.
\end{enumerate}
代数簇的范畴是代数空间范畴的全子范畴.代数空间比代数簇有许多优点，我只提两个.

代数空间在群作用下的商更可能存在.例如，Hironaka 构造了非奇异 3 维簇 $V$ 和 2 阶群 $G$ 在 $V$ 上的自由作用，使得商 $V/G$ 不作为代数簇存在(问题是存在不包含在任何仿射开集中的 $G$ 轨道).有限群自由作用下的代数空间的商总是作为代数空间存在(Artin 1973, p.109).关于代数空间商的非常一般的定理，见 Keel, S., and Mori, S., \emph{Quotients by groupoids}, Annals of Mathematics, 145 (1997), 193--213.

从 $\mathbb{C}$ 上完备代数簇的范畴到紧解析空间范畴有自然的全忠实函子 $V \mapsto V(\mathbb{C})$，但对这个函子的本质像没有方便的刻画.对完备代数空间范畴有：它由这样的紧解析空间 $M$ 组成：$M$ 上亚纯函数域有等于 $M$ 的维数的超越次数(这些是 Moishezon 空间).

\subsection*{附记}

设 $X$ 是簇(或拟紧概形).则 $X$ 的每个开子集 $U$ 都拟紧，即 $U$ 的开子集覆盖 $U = \bigcup U_i$ 都有有限子覆盖.此外，由于我把平展态射定义为有限型(而不是局部有限型)，每个在 $X$ 上平展的 $U$ 都拟紧.因为平展态射是开的，这蕴含每个平展覆盖 $(U_i \to U)_{i \in I}$ 都有有限子覆盖.每个覆盖族都含有限覆盖子族的景称为 Noether 景.注意拓扑空间 $X$ 定义的景很少是 Noether 的：若 $X$ 是 Hausdorff 的，它的景是 Noether 的，当且仅当 $X$ 的每个开子集都紧.

\begin{quote}
\refstepcounter{theorem}\label{item:7.20}
\textbf{7.20} 对 Noether 景，要证明预层是层，只需对有限覆盖检验层条件.由此推出，按朴素(即预层)方式形成的层的直和与正向极限仍是层，并且由此推出上同调与直和及(对有向集的)正向极限交换.见 EC III.3.
\end{quote}
